\begin{proof}[Proof of \cref{obj:lem:local-polynomial-projection-facts}]
All constants constructed below depend only on the fixed public parameters. Write
\[
\begin{gathered}
\mathcal U_{\mathrm{coarse}}
 =\left\{\kappa\in\mathbb N^d:|\kappa|\le2\right\},
\qquad
|\kappa|=\sum_{i=1}^d\kappa_i,
\qquad
\kappa!=\prod_{i=1}^d\kappa_i!,
\qquad
q=|\mathcal U_{\mathrm{coarse}}|,\\
\mathsf L_0(\upsilon)=1,\qquad
\mathsf L_1(\upsilon)=\sqrt{12}\,\upsilon,\qquad
\mathsf L_2(\upsilon)=\sqrt5\,(6\upsilon^2-1/2)
\quad(\upsilon\in\mathbb R).
\end{gathered}
\]
Thus the coarse coordinates are
\[
r_\kappa(x)
 =\prod_{i=1}^d
   \mathsf L_{\kappa_i}\!\left(\frac{x_i-1/2}{h}\right),
\qquad
p_\vartheta(x)=\sum_{\kappa\in\mathcal U_{\mathrm{coarse}}}
r_\kappa(x)\vartheta_\kappa
\quad(x\in\mathbb R^d).
\]
We use the designated functions supplied by \cref{obj:def:primitive-class} and the localized measure and projection of \cref{obj:def:projection}.

\begin{enumerate}
\item \emph{Taylor coefficients and approximation.}
We first construct a positive constant \(C_{\mathrm T}\), uniform over admitted laws and localization scales.

If \(\gamma=1\), then \(p_*=0\). Let
\[
\mathbf0=(0,\ldots,0),\qquad
\vartheta_\kappa=
\begin{cases}
\tau_P(x_0),&\kappa=\mathbf0,\\
0,&\kappa\ne\mathbf0,
\end{cases}
\qquad
C_{\mathrm T}=L(\sqrt d+1).
\]
Since \(r_{\mathbf0}=1\), this gives
\[
p_\vartheta(x)=\tau_P(x_0)=p(x).
\]
By \cref{obj:def:holder-norm,obj:def:primitive-class},
\[
|\tau_P(x_0)|\le L,
\qquad
|\tau_P(x)-\tau_P(x')|\le L\|x-x'\|_2
\quad(x,x'\in[0,1]^d).
\]
For equal arguments the second inequality follows from its zero left-hand side; for distinct arguments it is the order-zero Hölder bound. Moreover, for \(x\in C_h\),
\[
\|x-x_0\|_2^2
 =\sum_{i=1}^d|x_i-1/2|^2
 \le dh^2.
\]
The inclusion \(C_h\subset[0,1]^d\) follows from \(h\le1/2\). Consequently,
\[
|\tau_P(x)-p_\vartheta(x)|
 \le L\sqrt d\,h
 \le C_{\mathrm T}h,
\qquad
\|\vartheta\|_2=|\tau_P(x_0)|\le L\le C_{\mathrm T}.
\]
The representation also gives \(p_\vartheta(x_0)=\tau_P(x_0)\).

Suppose now that \(\gamma>1\). Define
\[
p_*=\lceil\gamma\rceil-1,\qquad
s_{\mathrm T}=\gamma-p_*,
\qquad
D_{\mathrm{dim}}=d+1.
\]
The public range of \(\gamma\) implies
\[
p_*\in\{1,2\},\qquad
0<s_{\mathrm T}\le1,\qquad
p_*+s_{\mathrm T}=\gamma.
\]
In particular, at \(\gamma=2\) and \(\gamma=3\), the pairs
\((p_*,s_{\mathrm T})\) are \((1,1)\) and \((2,1)\), respectively.
The same two definitions supply continuous derivatives through order \(p_*\) and the bounds
\[
\begin{aligned}
|\partial^\kappa\tau_P(x)|&\le L
&&(|\kappa|\le p_*,\ x\in[0,1]^d),\\
|\partial^\kappa\tau_P(x)-\partial^\kappa\tau_P(x')|
 &\le L\|x-x'\|_2^{s_{\mathrm T}}
&&(|\kappa|=p_*,\ x,x'\in[0,1]^d).
\end{aligned}
\]
The second inequality includes equal arguments, where both sides vanish.

For \(x\in C_h\), set
\[
z_x=x-x_0,\qquad
F_x(\upsilon)=\tau_P(x_0+\upsilon z_x)
\quad(0\le\upsilon\le1).
\]
The segment lies in the design cube. The chain rule, together with symmetry of the continuous derivatives, yields
\[
F_x^{(p_*)}(\upsilon)
 =
 \sum_{(i_1,\ldots,i_{p_*})\in\{1,\ldots,d\}^{p_*}}
 \left(\prod_{j=1}^{p_*}z_{x,i_j}\right)
 (\partial_{i_1}\cdots\partial_{i_{p_*}}\tau_P)
 (x_0+\upsilon z_x).
\]
Each coordinate of \(z_x\) has absolute value at most \(\|z_x\|_2\).
There are \(d^{p_*}\le D_{\mathrm{dim}}^{p_*}\) terms, so the top-derivative modulus gives
\[
\left|F_x^{(p_*)}(\upsilon)-F_x^{(p_*)}(\upsilon')\right|
 \le D_{\mathrm{dim}}^{p_*}L
       \|z_x\|_2^{p_*+s_{\mathrm T}}
       |\upsilon-\upsilon'|^{s_{\mathrm T}}
\quad(0\le\upsilon,\upsilon'<1).
\]
Grouping the coordinate derivatives by their multiplicities gives
\[
p(x)
 =\sum_{j=0}^{p_*}\frac{F_x^{(j)}(0)}{j!}
 =\sum_{|\kappa|\le p_*}
   \frac{\partial^\kappa\tau_P(x_0)}{\kappa!}
   \prod_{i=1}^d(x_i-1/2)^{\kappa_i}.
\]
The Lagrange remainder for the Taylor polynomial of degree \(p_*-1\)
provides \(\upsilon_x\in(0,1)\) such that, after subtracting the degree-\(p_*\) term,
\[
\tau_P(x)-p(x)
 =\frac{F_x^{(p_*)}(\upsilon_x)-F_x^{(p_*)}(0)}{p_*!}.
\]
Since \(p_*!\ge1\) and \(\upsilon_x^{s_{\mathrm T}}\le1\), this implies
\[
|\tau_P(x)-p(x)|
 \le D_{\mathrm{dim}}^{p_*}L\|z_x\|_2^\gamma.
\]
To obtain a constant uniform in the fractional order, observe that
\[
\|z_x\|_2\le\sqrt d\,\frac h2\le D_{\mathrm{dim}}h,
\qquad
D_{\mathrm{dim}}\ge1,
\qquad
\gamma=p_*+s_{\mathrm T}\le p_*+1.
\]
Hence, with
\[
C_{\mathrm{rem}}
 =D_{\mathrm{dim}}^{p_*}D_{\mathrm{dim}}^{p_*+1},
\]
we have
\[
|\tau_P(x)-p(x)|
 \le C_{\mathrm{rem}}Lh^\gamma
 \quad(x\in C_h).
\]

We next represent \(p\) in the coarse basis by an explicit change of coefficients.
Define the matrix, whose rows and columns are indexed by \(0,1,2\),
\[
\mathsf M_h=
\begin{pmatrix}
1&0&0\\
0&h/\sqrt{12}&0\\
h^2/12&0&h^2/(6\sqrt5)
\end{pmatrix}.
\]
Direct substitution gives
\[
(h\upsilon)^j
 =\sum_{\ell=0}^2(\mathsf M_h)_{j\ell}\mathsf L_\ell(\upsilon)
 \quad(j\in\{0,1,2\},\ \upsilon\in\mathbb R).
\]
Its entries satisfy
\[
(\mathsf M_h)_{j\ell}=0\quad(\ell>j),
\qquad
|(\mathsf M_h)_{j\ell}|\le h^j.
\]
Put
\[
\begin{gathered}
\mathcal U_{\mathrm T}
 =\{\kappa\in\mathbb N^d:|\kappa|\le p_*\},
\qquad
N_{\mathrm T}=|\mathcal U_{\mathrm T}|,\\
a_{\mathrm T,\kappa}
 =\frac{\partial^\kappa\tau_P(x_0)}{\kappa!}
 \quad(\kappa\in\mathcal U_{\mathrm T}),\\
\mathsf M_h[\kappa,\kappa']
 =\prod_{i=1}^d(\mathsf M_h)_{\kappa_i\kappa'_i}
 \quad(\kappa,\kappa'\in\mathcal U_{\mathrm T}).
\end{gathered}
\]
Since \(\kappa!\ge1\),
\[
|a_{\mathrm T,\kappa}|\le L,
\qquad
|\mathsf M_h[\kappa,\kappa']|
 \le h^{|\kappa|}\le1.
\]
Expanding the product of the one-dimensional identities gives
\[
\prod_{i=1}^d(x_i-1/2)^{\kappa_i}
 =\sum_{\kappa'\in\mathcal U_{\mathrm T}}
   \mathsf M_h[\kappa,\kappa']r_{\kappa'}(x).
\]
Indeed, a nonzero term has \(\kappa'_i\le\kappa_i\) in every coordinate,
so its total degree remains at most \(p_*\).

Define the coarse vector by zero extension:
\[
\vartheta_{\kappa'}=
\begin{cases}
\displaystyle\sum_{\kappa\in\mathcal U_{\mathrm T}}
 a_{\mathrm T,\kappa}\mathsf M_h[\kappa,\kappa'],
 &\kappa'\in\mathcal U_{\mathrm T},\\
0,&\kappa'\in\mathcal U_{\mathrm{coarse}}\setminus\mathcal U_{\mathrm T}.
\end{cases}
\]
Exchanging the finite sums proves
\[
p_\vartheta(x)=p(x)\quad(x\in\mathbb R^d).
\]
Furthermore,
\[
|\vartheta_{\kappa'}|\le N_{\mathrm T}L,
\qquad
\sum_{\kappa'\in\mathcal U_{\mathrm{coarse}}}
 \vartheta_{\kappa'}^2
 \le N_{\mathrm T}(N_{\mathrm T}L)^2.
\]
Thus, defining
\[
C_{\mathrm{coef}}=\sqrt{N_{\mathrm T}}\,N_{\mathrm T}+1,
\qquad
C_{\mathrm T}=L\max\{C_{\mathrm{rem}},C_{\mathrm{coef}}\},
\]
we obtain
\[
\|\vartheta\|_2\le C_{\mathrm{coef}}L\le C_{\mathrm T},
\qquad
\sup_{x\in C_h}|\tau_P(x)-p_\vartheta(x)|
 \le C_{\mathrm T}h^\gamma.
\]
At \(x_0\), every positive-degree centered monomial vanishes and the
zero-degree coefficient is \(\tau_P(x_0)\); therefore
\(p_\vartheta(x_0)=\tau_P(x_0)\).
Both cases construct a finite positive \(C_{\mathrm T}\).
% lean: CausalSmith.Stat.TwosamplePointcateAnnotationFrontier.taylor_coefficient_bounds

\item \emph{Basis and kernel bounds.}
For \(|\upsilon|\le1/2\), the three reference polynomials satisfy
\[
|\mathsf L_0(\upsilon)|=1,\qquad
|\mathsf L_1(\upsilon)|\le\frac{\sqrt{12}}2\le4,\qquad
|\mathsf L_2(\upsilon)|\le\sqrt5\le4.
\]
For the last inequality we used
\(0\le\upsilon^2\le1/4\), which gives
\(|6\upsilon^2-1/2|\le1\). It follows that
\[
|r_\kappa(x)|\le4^d,
\qquad
\|r(x)\|_2^2\le q(4^d)^2
\quad(x\in C_h).
\]

Here is the fine-cell apparatus used for the kernel calculations.
Define
\[
\delta=\frac hJ>0,\qquad
\mathcal Z_J=\{0,\ldots,J-1\}^d,
\qquad
c_{\mathbf j,i}=1/2-h/2+(j_i+1/2)\delta
\quad(\mathbf j\in\mathcal Z_J).
\]
For each coordinate, let
\[
I_{j_i}=
\begin{cases}
[1/2-h/2+j_i\delta,\;1/2-h/2+(j_i+1)\delta),
 &j_i<J-1,\\
[1/2-h/2+(J-1)\delta,\;1/2+h/2],
 &j_i=J-1,
\end{cases}
\qquad
D_{\mathbf j}=\prod_{i=1}^d I_{j_i}.
\]
These cells partition \(C_h\), with every terminal face assigned to the
last interval. To see coverage explicitly, for \(x\in C_h\) define
\[
t_{x,i}=\frac{x_i-(1/2-h/2)}{\delta}\in[0,J],
\qquad
j_i(x)=
\begin{cases}
J-1,&t_{x,i}=J,\\
\lfloor t_{x,i}\rfloor,&t_{x,i}<J.
\end{cases}
\]
Then \(x\in D_{\mathbf j(x)}\). Distinct indices give disjoint cells:
in a coordinate where their indices differ, the smaller interval's right
endpoint is at or before the larger interval's left endpoint, and the
smaller interval excludes its right endpoint.
These conventions also apply when \(J=1\).

The fine coordinates are
\[
b_{\mathbf j,\kappa}(x)=
\begin{cases}
\displaystyle
\sqrt{J^d}\prod_{i=1}^d
 \mathsf L_{\kappa_i}\!\left(\frac{x_i-c_{\mathbf j,i}}{\delta}\right),
 &x\in D_{\mathbf j},\\
0,&x\notin D_{\mathbf j},
\end{cases}
\quad
(\mathbf j\in\mathcal Z_J,\ \kappa\in\mathcal U_{\mathrm{coarse}}).
\]
In particular, all fine coordinates vanish outside \(C_h\).
The orthonormality specified in \cref{obj:def:projection} reads
\[
\int b_{\mathbf j,\kappa}(x)b_{\mathbf j',\kappa'}(x)\,d\nu_h(x)
 =
\begin{cases}
1,&(\mathbf j,\kappa)=(\mathbf j',\kappa'),\\
0,&(\mathbf j,\kappa)\ne(\mathbf j',\kappa').
\end{cases}
\]
Its normalization can also be seen directly: the reference integrals satisfy
\[
\int_{-1/2}^{1/2}\mathsf L_j(\upsilon)\mathsf L_\ell(\upsilon)\,d\upsilon
 =\mathbf1_{\{j=\ell\}}
 \quad(j,\ell\in\{0,1,2\}),
\]
by integrating the displayed polynomials. On a common cell, the
change of variables \(x_i=c_{\mathbf j,i}+\delta\upsilon_i\) contributes
\[
h^{-d}J^d\delta^d=1.
\]
Different cells have disjoint supports, and the assigned faces have zero
Lebesgue measure.

The coordinate bound gives, for every \(x\in\mathbb R^d\),
\[
|b_{\mathbf j,\kappa}(x)|\le\sqrt{J^d}\,4^d.
\]
The constant coordinate on a cell is
\[
b_{\mathbf j,\mathbf0}(x)=\sqrt{J^d}\,\mathbf1_{D_{\mathbf j}}(x).
\]
Consequently,
\[
|b_{\mathbf j,\kappa}(x)|
 \le\frac{4^d}{\sqrt{J^d}}b_{\mathbf j,\mathbf0}(x)^2,
\qquad
\int|b_{\mathbf j,\kappa}(x)|\,d\nu_h(x)
 \le\frac{4^d}{\sqrt{J^d}},
\]
where the second inequality uses
\(\int b_{\mathbf j,\mathbf0}^2\,d\nu_h=1\).

If \(x\in D_{\mathbf j}\), disjointness reduces the kernel to
\[
K_J(x,x')
 =\sum_{\kappa\in\mathcal U_{\mathrm{coarse}}}
 b_{\mathbf j,\kappa}(x)b_{\mathbf j,\kappa}(x').
\]
The triangle inequality and the preceding bounds therefore give
\[
\begin{aligned}
\int|K_J(x,x')|\,d\nu_h(x')
&\le
 \sum_{\kappa\in\mathcal U_{\mathrm{coarse}}}
 |b_{\mathbf j,\kappa}(x)|
 \int|b_{\mathbf j,\kappa}(x')|\,d\nu_h(x')\\
&\le q\bigl(\sqrt{J^d}\,4^d\bigr)
          \frac{4^d}{\sqrt{J^d}}
 =q(4^d)^2.
\end{aligned}
\]
For \(x\) outside all cells, \(K_J(x,x')=0\), so the same bound holds there.

There are \(|\mathcal Z_J|=J^d\) cells and \(q\) coordinates per cell;
hence the fine index set has \(qJ^d\) elements. Expanding the squared
finite kernel and factoring the product integrals yields
\[
\begin{aligned}
\iint K_J(x,x')^2\,d\nu_h(x)\,d\nu_h(x')
&=
 \sum_{\mathbf j,\kappa}\sum_{\mathbf j',\kappa'}
 \left(\int
 b_{\mathbf j,\kappa}(x)b_{\mathbf j',\kappa'}(x)\,d\nu_h(x)\right)^2\\
&=qJ^d=k.
\end{aligned}
\]
All terms are integrable by the square integrability of the basis coordinates.
Define
\[
C_{\mathrm B}
 =\max\{\sqrt q\,4^d+1,\ q(4^d)^2+1\}>0.
\]
We have proved
\[
\sup_{x\in C_h}\|r(x)\|_2\le C_{\mathrm B},
\qquad
\sup_{x\in C_h}\int|K_J(x,x')|\,d\nu_h(x')\le C_{\mathrm B},
\]
together with the exact squared-kernel identity.
% lean: CausalSmith.Stat.TwosamplePointcateAnnotationFrontier.kernel_basis_bounds

\item \emph{Control-mean residual.}
For square-integrable functions, write
\[
\|f\|_h=\left(\int f(x)^2\,d\nu_h(x)\right)^{1/2}.
\]
The coarse constant coordinate equals one, so coarse orthonormality gives
\[
\nu_h(C_h)=\int1\,d\nu_h
 =\int r_{\mathbf0}(x)^2\,d\nu_h(x)=1.
\]
Thus \(\nu_h\) is a probability measure supported on \(C_h\).

The centers satisfy \(c_{\mathbf j}\in D_{\mathbf j}\).
For two points in the same cell, coordinatewise interval bounds give
\[
|x_i-x_i'|\le\delta,
\qquad
\|x-x'\|_2^2\le d\delta^2
\quad(x,x'\in D_{\mathbf j}).
\]
By \cref{obj:def:holder-norm,obj:def:primitive-class} and
\(0<\beta\le1\),
\[
|\mu_{0,P}(x)|\le L,
\qquad
|\mu_{0,P}(x)-\mu_{0,P}(x')|
 \le L\|x-x'\|_2^\beta
\quad(x,x'\in C_h).
\]
Indeed, \(\lceil\beta\rceil-1=0\), so the Hölder norm bounds the
function and its order-zero modulus; equal arguments satisfy the modulus
inequality directly.

Define the fine coefficient vector and its expansion by
\[
v^{(0,J)}_{\mathbf j,\kappa}=
\begin{cases}
\mu_{0,P}(c_{\mathbf j})/\sqrt{J^d},&\kappa=\mathbf0,\\
0,&\kappa\ne\mathbf0,
\end{cases}
\qquad
f_{0,J}(x)=
 \sum_{\mathbf j,\kappa}
 b_{\mathbf j,\kappa}(x)v^{(0,J)}_{\mathbf j,\kappa}.
\]
The cell assignments imply
\[
f_{0,J}(x)=
\begin{cases}
\mu_{0,P}(c_{\mathbf j}),&x\in D_{\mathbf j},\\
0,&x\notin C_h.
\end{cases}
\]
Set
\[
E_{0,J}=L(\sqrt d)^\beta\delta^\beta\ge0.
\]
For \(x\in D_{\mathbf j}\), the modulus and distance bounds give
\[
|\mu_{0,P}(x)-f_{0,J}(x)|
 \le L\|x-c_{\mathbf j}\|_2^\beta
 \le L(\sqrt d\,\delta)^\beta
 =E_{0,J}.
\]
Because \(\nu_h\) is a probability measure,
\[
\|\mu_{0,P}-f_{0,J}\|_h^2
 \le\int E_{0,J}^2\,d\nu_h=E_{0,J}^2.
\]

We now derive the projection comparison used here and in the next step.
For any \(f\in L_2(\nu_h)\) and any real fine coefficient vector
\(v_J=(v_{J,\mathbf j,\kappa})\), put
\[
f_J(x)=\sum_{\mathbf j,\kappa}
 b_{\mathbf j,\kappa}(x)v_{J,\mathbf j,\kappa}.
\]
Expanding the kernel integral in \cref{obj:def:projection} gives
\[
\Pi_Jf(x)=
 \sum_{\mathbf j,\kappa}b_{\mathbf j,\kappa}(x)
 \int b_{\mathbf j,\kappa}(x')f(x')\,d\nu_h(x').
\]
Orthonormality then implies
\[
\int b_{\mathbf j,\kappa}(x)f_J(x)\,d\nu_h(x)
 =v_{J,\mathbf j,\kappa},
\qquad
\Pi_J(f-f_J)=\Pi_Jf-f_J,
\]
and
\[
\int b_{\mathbf j,\kappa}(x)
       \bigl(f(x)-\Pi_Jf(x)\bigr)\,d\nu_h(x)=0.
\]
The last identity makes the cross term in the squared norm vanish, giving
\[
\|f\|_h^2=\|\Pi_Jf\|_h^2+\|f-\Pi_Jf\|_h^2.
\]
Apply this identity to \(f-f_J\) and use the displayed projection identity:
\[
\|f-f_J\|_h^2
 =\|\Pi_Jf-f_J\|_h^2+\|f-\Pi_Jf\|_h^2
 \ge\|f-\Pi_Jf\|_h^2.
\]
The bounded control mean and its finite-basis approximant are square
integrable. Taking \(f=\mu_{0,P}\) and \(f_J=f_{0,J}\), and then taking
square roots, proves
\[
\|\mu_{0,P}-\Pi_J\mu_{0,P}\|_h
 \le E_{0,J}.
\]
Therefore the positive constant
\[
C_{\mathrm m}=L(\sqrt d)^\beta+1
\]
satisfies
\[
\|\mu_{0,P}-\Pi_J\mu_{0,P}\|_h
 \le C_{\mathrm m}\delta^\beta.
\]
% lean: CausalSmith.Stat.TwosamplePointcateAnnotationFrontier.control_projection_bounds

\item \emph{Propensity-weighted residuals.}
Similarly, \(0<\alpha\le1\) and
\cref{obj:def:holder-norm,obj:def:primitive-class} give
\[
|e_P(x)|\le L,
\qquad
|e_P(x)-e_P(x')|\le L\|x-x'\|_2^\alpha
\quad(x,x'\in C_h).
\]
Fix \(u\in\mathcal U_{\mathrm{coarse}}\), and define
\[
f_{e,u}(x)=e_P(x)r_u(x)
\quad(x\in\mathbb R^d).
\]
On \(C_h\), \(|f_{e,u}(x)|\le L4^d\), so it is square integrable
under \(\nu_h\).

We construct a fine-basis expansion that equals
\(e_P(c_{\mathbf j})r_u\) on each cell. Define
\[
\zeta_{\mathbf j,i}=\frac{c_{\mathbf j,i}-1/2}{h},
\qquad
\sigma_J=\frac1J,
\qquad
\mathsf A_{\mathbf j,i}=
\begin{pmatrix}
1&0&0\\
\sqrt{12}\,\zeta_{\mathbf j,i}&\sigma_J&0\\
\sqrt5(6\zeta_{\mathbf j,i}^2+\sigma_J^2/2-1/2)
 &\sqrt5\sqrt{12}\,\zeta_{\mathbf j,i}\sigma_J&\sigma_J^2
\end{pmatrix}.
\]
Rows and columns are indexed by \(0,1,2\). Substitution into the three
reference polynomials verifies
\[
\mathsf L_j(\zeta_{\mathbf j,i}+\sigma_J\upsilon_{\mathrm{aff}})
 =
 \sum_{\ell=0}^j
 (\mathsf A_{\mathbf j,i})_{j\ell}
 \mathsf L_\ell(\upsilon_{\mathrm{aff}})
\quad(j\in\{0,1,2\},\ \upsilon_{\mathrm{aff}}\in\mathbb R).
\]
For \(x\in D_{\mathbf j}\),
\[
\frac{x_i-1/2}{h}
 =\zeta_{\mathbf j,i}
   +\sigma_J\frac{x_i-c_{\mathbf j,i}}{\delta}.
\]
Define the tensor coefficients
\[
\mathsf A_{\mathbf j}[u,\kappa]
 =\prod_{i=1}^d(\mathsf A_{\mathbf j,i})_{u_i\kappa_i}
 \quad(\kappa\in\mathcal U_{\mathrm{coarse}}).
\]
Expanding the product of the affine identities gives
\[
r_u(x)=
 \sum_{\kappa\in\mathcal U_{\mathrm{coarse}}}
 \mathsf A_{\mathbf j}[u,\kappa]
 \prod_{i=1}^d
 \mathsf L_{\kappa_i}\!\left(\frac{x_i-c_{\mathbf j,i}}{\delta}\right).
\]
Only indices with \(\kappa_i\le u_i\) contribute, and hence every
contributing index has total degree at most two.

Now set
\[
v^{(e,u,J)}_{\mathbf j,\kappa}
 =\frac{e_P(c_{\mathbf j})\mathsf A_{\mathbf j}[u,\kappa]}
        {\sqrt{J^d}},
\qquad
f_{e,u,J}(x)=
 \sum_{\mathbf j,\kappa}
 b_{\mathbf j,\kappa}(x)v^{(e,u,J)}_{\mathbf j,\kappa}.
\]
Disjointness of the cells and the tensor identity imply the everywhere
defined formula
\[
f_{e,u,J}(x)=
\begin{cases}
e_P(c_{\mathbf j})r_u(x),&x\in D_{\mathbf j},\\
0,&x\notin C_h.
\end{cases}
\]
Define
\[
E_{e,J}=L(\sqrt d)^\alpha4^d\delta^\alpha\ge0.
\]
For \(x\in D_{\mathbf j}\),
\[
\begin{aligned}
|f_{e,u}(x)-f_{e,u,J}(x)|
 &=|e_P(x)-e_P(c_{\mathbf j})|\,|r_u(x)|\\
 &\le L\|x-c_{\mathbf j}\|_2^\alpha4^d\\
 &\le L(\sqrt d\,\delta)^\alpha4^d
 =E_{e,J}.
\end{aligned}
\]
The probability normalization therefore gives
\[
\|f_{e,u}-f_{e,u,J}\|_h^2\le E_{e,J}^2.
\]
Applying the projection comparison from the preceding step to this
fine-basis expansion yields
\[
\|e_Pr_u-\Pi_J(e_Pr_u)\|_h
 \le\|f_{e,u}-f_{e,u,J}\|_h
 \le E_{e,J}.
\]
Thus the positive constant
\[
C_{\mathrm e}=L(\sqrt d)^\alpha4^d+1
\]
satisfies, uniformly in \(u\),
\[
\|e_Pr_u-\Pi_J(e_Pr_u)\|_h
 \le C_{\mathrm e}\delta^\alpha.
\]
% lean: CausalSmith.Stat.TwosamplePointcateAnnotationFrontier.propensity_projection_bounds

\item \emph{A common residual constant.}
Define
\[
C_{\mathrm R}=\max\{C_{\mathrm e},C_{\mathrm m}\}>0.
\]
Since \(\delta^\alpha,\delta^\beta\ge0\), the preceding estimates imply
\[
\|e_Pr_u-\Pi_J(e_Pr_u)\|_h
 \le C_{\mathrm R}\delta^\alpha
 \quad(u\in\mathcal U_{\mathrm{coarse}}),
\qquad
\|\mu_{0,P}-\Pi_J\mu_{0,P}\|_h
 \le C_{\mathrm R}\delta^\beta.
\]
% lean: CausalSmith.Stat.TwosamplePointcateAnnotationFrontier.residual_projection_bounds

\item \emph{Combination of the bounds.}
Finally, choose
\[
C=\max\{C_{\mathrm T},\max\{C_{\mathrm B},C_{\mathrm R}\}\}.
\]
This is finite, positive, and depends only on the public parameters.
The inequalities
\[
C_{\mathrm T}\le C,\qquad
C_{\mathrm B}\le C,\qquad
C_{\mathrm R}\le C,
\qquad
h^\gamma\ge0,\qquad
\delta^\alpha\ge0,\qquad
\delta^\beta\ge0
\]
allow every bound above to be enlarged to the claimed constant \(C\).
The same Taylor coefficient vector retains its exact center value and
its equality with \(p\) on \(C_h\), while the squared-kernel integral
retains the exact value \(qJ^d=k\). These conclusions hold for every
admitted \(P\), every \(0<h\le1/2\), and every integer \(J\ge1\).
% lean: CausalSmith.Stat.TwosamplePointcateAnnotationFrontier.local_polynomial_projection_facts
\end{enumerate}
\end{proof}
